\documentclass[12pt]{article}
\usepackage[margin=1in]{geometry}
\setlength{\parskip}{0.5em}

\title{Reference Items Across Several Domains}
\author{legible-pdf synthesized corpus}
\date{}

\begin{document}
\maketitle

\section{Inner Planets}

The inner solar system contains rocky bodies orbiting closest to the Sun.

\begin{itemize}
\item Mercury
\item Venus
\item Earth
\end{itemize}

Astronomers study these bodies using ground-based and orbital telescopes.

\section{Primary Colors}

Three colors form the additive primary set used in display technology.

\begin{enumerate}
\item Red
\item Green
\item Blue
\end{enumerate}

These wavelengths combine to produce the visible spectrum on emissive screens.

\section{Produce Categories}

Markets organize fresh produce into broad categories with member items.

\begin{itemize}
\item Fruits
  \begin{itemize}
  \item Apple
  \item Banana
  \end{itemize}
\item Vegetables
  \begin{itemize}
  \item Carrot
  \item Daikon
  \end{itemize}
\end{itemize}

These categories shape how grocers arrange shelves and how shoppers plan meals.

\section{Programming Language Terms}

Compiler theory distinguishes several program transformation roles.

\begin{description}
\item[Compiler] A program that translates source code into executable code.
\item[Interpreter] A program that executes source code line by line.
\item[Linker] A program that combines compiled object files into an executable.
\end{description}

These roles often overlap in modern toolchains but remain conceptually distinct.

\section{Glossary}

The document mentions: Mercury, Venus, Earth, Red, Green, Blue, Fruits,
Apple, Banana, Vegetables, Carrot, Daikon, Compiler, Interpreter, Linker.

\end{document}
